\chapter{Introducción}
\label{sec:paper-intro}

\section{Motivación}

Un analista que investiga un incidente de seguridad raramente se
hace preguntas tabulares. Las preguntas tienen forma relacional:
\emph{¿qué usuarios autenticaron desde la IP que conectó al host
comprometido, dentro de una ventana de una hora?}. El SIEM entiende
\emph{usuario}, \emph{IP} y \emph{host} como columnas y la
pregunta, como una serie de joins textuales. La estructura
relacional del incidente vive en la cabeza del analista.
GraphHunter la explicita y la mueve al código.

\section{Problema}

Formalizamos el problema en tres subproblemas con propiedades
distintas:

\begin{enumerate}
    \item \textbf{Ingesta heterogénea}. Datos provienen de fuentes
    con esquemas distintos (EDR, firewalls, proxies, IAM, cloud
    audit). Los formatos mutan en el tiempo. El objetivo es un
    grafo canonical homogéneo sin perder la fidelidad del origen.
    \item \textbf{Matching estructural}. Dado un grafo
    $G = (V, E)$ con nodos y aristas tipados y una hipótesis
    (patrón en un DSL), enumerar \emph{paths} instancia del patrón
    con cotas temporales y estructurales. El espacio es
    combinatorio: naïf explota con la multiplicación de vecindarios.
    \item \textbf{Priorización}. El \emph{matching} produce muchos
    paths; la mayoría son actividad legítima. Un analista puede
    triar orden de decenas por turno. El objetivo es un score
    compuesto, interpretable y con pesos operables.
\end{enumerate}

\section{Alcance y no-alcance}

Este reporte describe las decisiones de diseño que fijan la forma
del sistema actual (\emph{post-shearing-layers}) y el razonamiento
detrás de cada una, con evaluación cuantitativa donde hay números y
honesta donde no. No es un paper académico con peer review; es un
reporte interno para consumo del equipo y eventualmente de
auditores técnicos externos.

Quedan fuera de alcance: el motor de UI
(\code{apps/tauri}), las políticas organizacionales que usan
GraphHunter como herramienta de triage, y la operación del
\emph{stack} de observabilidad sobre el que corre. Esos temas
tienen su propio lugar en el \emph{manual} y en la
\emph{especificación}.

\section{Contribuciones}

Este trabajo aporta:

\begin{enumerate}
    \item Una descomposición en \emph{shearing layers}
    (cap.~\ref{sec:paper-arq}) que reduce acoplamiento y permite
    evolucionar superficies sin tocar el núcleo.
    \item Un matcher con intersección SIMD de \emph{posting-lists}
    AVX2 (cap.~\ref{sec:paper-simd}) con equivalencia
    property-tested a una referencia DFS.
    \item Un DFS temporal guiado por score para el
    \emph{refinement} acotado por $k$-simplicity
    (cap.~\ref{sec:paper-dfs}).
    \item Un composer de anomalía de cinco componentes con pesos
    operables (cap.~\ref{sec:paper-score}).
    \item Un aislamiento FFI mediante \emph{typestate} + RAII +
    ABI versionada (cap.~\ref{sec:paper-ffi}).
    \item Un \emph{model gate} que vuelve opcional al clasificador
    GNN (cap.~\ref{sec:paper-gate}).
    \item Un ingestor de tres carriles con LLM-como-compilador
    offline y drift-como-señal (cap.~\ref{sec:paper-ingestor}).
\end{enumerate}

\section{Organización}

El capítulo~\ref{sec:paper-related} revisa enfoques previos.
El~\ref{sec:paper-arq} resume la arquitectura global. Los
capítulos~\ref{sec:paper-simd}--\ref{sec:paper-gate} desarrollan las
contribuciones técnicas. El~\ref{sec:paper-ingestor} cubre el
ingestor agéntico. El~\ref{sec:paper-eval} presenta evaluación. El
\ref{sec:paper-limites} discute limitaciones abiertas y el
\ref{sec:paper-conc} concluye.
